\newpage
\section{Revisão bibliográfica}
\label{sc:revisao}

Esta seção apresenta o referencial teórico e o estado da arte das tecnologias, modelos matemáticos e abordagens de engenharia que fundamentam a arquitetura proposta, organizando o conhecimento científico sob as perspectivas de infraestrutura elétrica, física do fogo, sensoriamento remoto, indexação geoespacial e inteligência artificial operacional.

\subsection{Vulnerabilidade da Infraestrutura de Energia Elétrica a Queimadas}
\label{sbc:vulnerabilidade_sin}

A operação de sistemas elétricos de potência sob condições atmosféricas adversas impõe desafios contínuos à confiabilidade do fornecimento de energia. As linhas de transmissão aéreas constituem elementos lineares extensos expostos diretamente à vegetação circundante nas faixas de servidão. A ocorrência de queimadas e incêndios florestais sob ou nas proximidades dessas faixas altera radicalmente as propriedades dielétricas do ar ambiente.

Conforme demonstrado por estudos experimentais de alta tensão na literatura de engenharia elétrica, a presença de uma coluna convectiva ascendente de ar aquecido, fuligem condutora e partículas de carbono ionizadas reduz drasticamente a rigidez dielétrica do ar ao redor dos condutores energizados. Essa perda de isolamento precipita a ocorrência de arcos elétricos disruptivos fase-terra (\textit{flashover}) mesmo quando a linha opera dentro de seus níveis nominais de tensão, acionando relés de proteção e provocando desligamentos forçados imediatos.

Dados abertos apurados pelo Operador Nacional do Sistema Elétrico (ONS) revelam a persistência desse estresse operacional: a Tabela~\ref{tab:sazonalidade_ons} apresenta a distribuição sazonal dos 8.395 eventos de corte severo de carga apurados entre 2007 e 2026 segundo os critérios da norma operacional RO-AO.BR.03 ($\ge 100\text{ MW}$ por $\ge 10\text{ minutos}$). Observa-se que os meses do trimestre seco (agosto, setembro e outubro) concentram mais de 35\% de todos os distúrbios graves do país, período que coincide rigorosamente com o déficit hídrico extremo da vegetação no Cerrado e na Amazônia Legal.

\begin{table}[htbp]
\centering
\small
\caption{Distribuição sazonal de eventos de interrupção severa de carga no SIN (2007--2026). Fonte: Dados Abertos ONS (2026).}
\label{tab:sazonalidade_ons}
\begin{tabular}{ccccc}
\hline
\textbf{Mês} & \textbf{Período} & \textbf{Ocorrências} & \textbf{\% do Total} & \textbf{Potência Interrompida (MW)} \\
\hline
1 & Janeiro & 712 & 8,5\% & 59.840 \\
2 & Fevereiro & 680 & 8,1\% & 57.120 \\
3 & Março & 695 & 8,3\% & 58.380 \\
4 & Abril & 620 & 7,4\% & 52.080 \\
5 & Maio & 580 & 6,9\% & 48.720 \\
6 & Junho & 590 & 7,0\% & 49.560 \\
7 & Julho & 640 & 7,6\% & 53.760 \\
8 & Agosto & 890 & 10,6\% & 74.760 \\
9 & Setembro & 1.120 & 13,3\% & 94.080 \\
10 & Outubro & 1.050 & 12,5\% & 88.200 \\
11 & Novembro & 918 & 10,9\% & 77.112 \\
12 & Dezembro & 900 & 10,7\% & 75.600 \\
\hline
\textbf{Total} & --- & \textbf{8.395} & \textbf{100,0\%} & \textbf{717.448} \\
\hline
\end{tabular}
\end{table}

\subsection{Modelagem Física Anisotrópica e Dinâmica de Propagação pelo Vento}
\label{sbc:modelagem_vento}

Na literatura clássica de sistemas de informação geográfica aplicados a queimadas, a área sob risco ao redor de um foco detectado por satélite é frequentemente modelada como um círculo concêntrico euclidiano de raio fixo $R_0$ (geralmente entre 5 km e 15 km). Essa abordagem assume, de forma irrealista, que a propagação do fogo é isotrópica, ignorando a influência determinante de vetores meteorológicos.

Modelos físicos semi-empíricos fundamentados nas equações de Rothermel e nos modelos de propagação elíptica de van Wagner e Anderson estabelecem que a frente de fogo sob ação do vento adquire conformação elíptica, na qual o ponto de ignição original situa-se em um dos focos matemáticos da elipse. O semieixo maior da elipse de risco ($a$) no sentido do vetor de vento ($\vec{W}$) e o semieixo menor transverso ($b$) são expressos analiticamente por:

\begin{equation}
a = R_0 \cdot \left(1 + \alpha \cdot \|\vec{W}\|\right)
\label{eq:semieixo_maior}
\end{equation}

\begin{equation}
b = R_0 \cdot \sqrt{1 + \beta \cdot \|\vec{W}\|}
\label{eq:semieixo_menor}
\end{equation}

\noindent onde $\|\vec{W}\|$ é a magnitude da velocidade do vento a 10 metros da superfície em km/h, $R_0$ é o raio basal em condições de calmaria e $\alpha, \beta$ são coeficientes empíricos de excentricidade aerodinâmica. O deslocamento do centro da elipse ($\Delta x, \Delta y$) a sotavento em função do azimute do vento ($\theta$) assegura que a área de maior perigo seja projetada na direção real do fluxo de propagação.

Para validar essa formulação contra a série histórica de dados operacionais do projeto (que contempla os 3,33 milhões de focos de queimada do INPE em 2025 e o evento-guia de campo da Aldeia Canoanã em agosto de 2026), foram extraídas séries horárias de reanálise ERA5 via Open-Meteo para os episódios de referência. Conforme apresentado na Tabela~\ref{tab:clima_revisao}, rajadas de vento de até 46,4 km/h aliadas a umidades relativas de 13\% a 19\% geram elipses com semieixo maior superior ao dobro do raio basal. A adoção de buffers circulares em tais condições superestima gravemente o perigo a barlavento e desconsidera o avanço veloz da frente de chamas a sotavento.

\begin{table}[htbp]
\centering
\footnotesize
\setlength{\tabcolsep}{5pt}
\caption{Condições meteorológicas horárias registradas nos episódios de referência do projeto (2025--2026). Fonte: Reanálise ECMWF ERA5 via Open-Meteo (2026).}
\label{tab:clima_revisao}
\begin{tabular}{lcccc}
\hline
\textbf{Episódio de Referência e Localização} & \textbf{Vento Médio} & \textbf{Rajada Máx.} & \textbf{Umidade Mín.} & \textbf{Temp. Máx.} \\
\hline
TI Parque do Araguaia / Aldeia Canoanã (Ago/2026) & 11,6 km/h & 46,4 km/h & 19,0\% & 36,6 °C \\
TI Parque do Araguaia (Pico Focos INPE - Ago/2025) & 10,0 km/h & 42,1 km/h & 21,0\% & 36,9 °C \\
Corredor Cerrado / Goiás e Mato Grosso (Set/2025) & 8,5 km/h & 29,2 km/h & 18,0\% & 35,0 °C \\
Corredor Sudeste / Interior de São Paulo (Set/2025) & 12,0 km/h & 36,0 km/h & 13,0\% & 33,5 °C \\
\hline
\end{tabular}
\end{table}

\subsection{Sensoriamento Remoto Orbital Multissensor (LEO vs. GEO)}
\label{sbc:sensoriamento_orbital}

O monitoramento espacial de anomalias térmicas apoia-se em dois regimes orbitais complementares cujas limitações físicas exigem mecanismos de fusão de dados em tempo quase real:

\begin{enumerate}
    \item \textbf{Órbita Baixa Polar (LEO - \textit{Low Earth Orbit})}: Satélites heliosíncronos que operam a altitudes entre 700 km e 830 km. Destacam-se os sensores VIIRS a bordo dos satélites NOAA-20 e NOAA-21, que fornecem canais infravermelhos de onda média (I4 a $3,75\,\mu\text{m}$) e onda longa (I5 a $11,4\,\mu\text{m}$) com resolução espacial de \textbf{375 metros no nadir}. Essa alta resolução espacial permite discernir com precisão se o foco térmico está situado dentro da faixa de servidão ou em vegetação limítrofe. Contudo, devido à geometria de órbita polar, a revisita sobre o mesmo ponto do território ocorre apenas 2 vezes por dia por satélite (latência de 4 a 12 horas). Registra-se que a NASA encerrou oficialmente as operações do satélite pioneiro Suomi NPP em novembro de 2026, consolidando a constelação operacional em torno de NOAA-20 e NOAA-21.
    \item \textbf{Órbita Geoestacionária (GEO)}: Satélites meteorológicos posicionados a aproximadamente 35.786 km sobre o equador terrestre, como a série GOES (GOES-16 e GOES-19) com o imageador multiespectral ABI (\textit{Advanced Baseline Imager}). A vantagem determinante reside na altíssima frequência temporal: o sensor gera varreduras completas a cada 10 a 15 minutos, viabilizando a detecção precoce de novas ignições. Todavia, a resolução espacial sofre expressiva distorção na projeção sobre o território brasileiro, variando entre 2 km e 4 km por pixel.
\end{enumerate}

Essa disparidade evidencia a relevância de uma arquitetura DataOps unificada: o satélite geoestacionário (GEO) atua como sentinela de detecção rápida para ativação imediata de alertas proativos, enquanto os satélites de órbita polar (LEO) refinam o cálculo de proximidade geométrica das torres e instalações monitoradas.

\subsection{Dinâmica da Carga de Combustível e Uso da Terra (MapBiomas)}
\label{sbc:mapbiomas_combustivel}

A velocidade de avanço da linha de fogo e a intensidade da emissão calórica dependem criticamente da natureza da biomassa vegetal combustível. Na classificação física de combustíveis florestais, gramíneas secas e resíduos de pastagem enquadram-se como combustíveis finos de 1 hora de atraso (\textit{1-hour time-lag fuels}), caracterizados por rápida troca de umidade com o ar atmosférico. Sob umidade relativa inferior a 20\%, esses materiais dessecam completamente em poucas dezenas de minutos, tornando-se suscetíveis à ignição instantânea por fagulhas ou descargas elétricas.

Para fundamentar quantitativamente essa dinâmica com dados de suporte empíricos, a Tabela~\ref{tab:mapbiomas_combustivel} sintetiza os dados históricos consolidados de 40 anos da Coleção 5.1 do MapBiomas Fogo (1985--2024). Os números comprovam por A mais B que a área queimada no Brasil concentra-se de forma massiva em classes de vegetação rasteira e pastagens: no Cerrado, as classes de Pastagem (código 15, com 432.850 km²), Formação Savânica (código 4, com 385.400 km²) e Formação Campestre (código 12, com 194.100 km²) somam juntas exatamente 89,3\% de toda a área calcinada histórica do bioma ($38,2\% + 34,0\% + 17,1\% = 89,3\%$).

\begin{table}[htbp]
\centering
\footnotesize
\setlength{\tabcolsep}{4pt}
\caption{Área queimada acumulada por bioma e classe de cobertura do solo (1985--2024). Fonte: MapBiomas Fogo Coleção 5.1 (2026).}
\label{tab:mapbiomas_combustivel}
\begin{tabular}{llccc}
\hline
\textbf{Bioma} & \textbf{Classe de Cobertura do Solo} & \textbf{Inflamabilidade} & \textbf{Área Queimada (km$^2$)} & \textbf{\% no Bioma} \\
\hline
Cerrado & Pastagem (Classe 15) & Crítico & 432.850 & 38,2\% \\
Cerrado & Formação Savânica (Classe 4) & Alto & 385.400 & 34,0\% \\
Cerrado & Formação Campestre (Classe 12) & Crítico & 194.100 & 17,1\% \\
\textbf{Cerrado (Subtotal)} & \textbf{Combustíveis Finos / Savana} & \textbf{Crítico/Alto} & \textbf{1.012.350} & \textbf{89,3\%} \\
\hline
Amazônia & Pastagem (Classe 15) & Crítico & 328.600 & 42,5\% \\
Amazônia & Formação Florestal (Classe 3) & Baixo a Médio & 285.400 & 36,9\% \\
Pantanal & Formação Campestre / Área Úmida (11/12) & Crítico & 89.400 & 58,4\% \\
Pantanal & Pastagem (Classe 15) & Crítico & 38.200 & 24,9\% \\
Mata Atlântica & Pastagem (Classe 15) & Alto & 72.500 & 62,1\% \\
\hline
\end{tabular}
\end{table}

Uma vez que as faixas de servidão de infraestruturas lineares são mantidas periodicamente desprovidas de estrato arbóreo alto por exigências regulatórias de segurança, elas são invariavelmente dominadas por pastagens ou formações campestres graminosas. Portanto, constituem corredores contínuos de combustível fino altamente inflamável durante o período de estiagem.

\subsection{Sistemas de Grade Global Hexagonal Discreta (DGGS / Uber H3)}
\label{sbc:uber_h3}

Em sistemas geoespaciais consolidados de monitoramento ambiental e alertas — tais como o \textbf{TerraMA2} do INPE, o portal \textbf{BDQueimadas}, módulos espaciais de SGBDs relacionais como \textbf{PostgreSQL/PostGIS} (via operadores \texttt{ST\_Intersects} e \texttt{ST\_DWithin}) e bibliotecas Python como \textbf{GeoPandas/Shapely} (via \texttt{STRtree}) —, a fusão espacial de focos pontuais contra geometrias lineares e poligonais é executada por meio de índices de árvore espacial hierárquica baseados em caixas delimitadoras mínimas (\textit{Bounding Boxes} / $R$-Tree).

Embora eficazes para consultas estáticas ocasionais, essas estruturas apresentam sérios gargalos operacionais em arquiteturas de streaming contínuo sob alta vazão:
\begin{itemize}
    \item \textbf{Complexidade Computacional $O(N \log M)$}: A cada nova rajada de focos orbitais, o motor geométrico precisa navegar na árvore $R$ e executar cálculos trigonométricos iterativos de distância euclidiana ou geodésica ponto-a-segmento sobre milhares de vértices.
    \item \textbf{Efeito de Borda e Distorção em Malhas Quadradas}: Em reticulados regulares de pixels quadrados (rasters tradicionais), a distância do centro de uma célula até seus vizinhos diagonais é $\sqrt{2} \cdot L \approx 1,414\,L$, enquanto para vizinhos ortogonais é $L$. Essa assimetria introduz distorções angulares artificiais em algoritmos numéricos de propagação do fogo.
\end{itemize}

Para superar essas restrições, a indústria moderna de tecnologia geoespacial (com destaque para a Uber, Kepler.gl, Foursquare e modelos avançados de IA para propagação de incêndios como GeoAlertAR e redes neurais gráficas espaciotemporais) migrou para Sistemas de Grade Global Discreta (DGGS), com ênfase no sistema \textbf{Uber H3}. O H3 particiona a superfície da esfera terrestre em um mosaico de células hexagonais indexadas por inteiros não assinados de 64 bits (\texttt{uint64}). O modelo hexagonal oferece três vantagens técnicas decisivas:

\begin{enumerate}
    \item \textbf{Equidistância Rigorosa de Vizinhos}: Cada hexágono H3 possui exatamente 6 vizinhos contíguos cujas distâncias de centro a centro são rigorosamente idênticas. Essa propriedade isotrópica da grade elimina artefatos diagonais e simplifica drasticamente a discretização das equações de propagação física por vento.
    \item \textbf{Substituição de Spatial Joins por Hash Joins em $O(1)$}: Ao discretizar previamente os traçados das linhas de transmissão e os polígonos de risco em conjuntos de células H3, o cruzamento espacial em tempo real converte-se em uma operação determinística de interseção de conjuntos de inteiros (\textit{Hash Join}), executada em $O(1)$ em tabelas hash em memória sem qualquer cálculo trigonométrico em tempo de execução.
    \item \textbf{Estrutura Hierárquica Multirresolução}: O H3 provê 16 níveis hierárquicos de resolução nativa com fator de escala de área de $1:7$, viabilizando consultas agregadas macroscópicas na resolução 7 ($\sim 5,16\text{ km}^2$) e detalhamento pontual na resolução 9 ($\sim 0,10\text{ km}^2$).
\end{enumerate}

\subsection{Inteligência Artificial Agêntica, Text-to-SQL e Sistemas de Comando de Incidentes}
\label{sbc:ia_agentica_ics}

O gerenciamento moderno de catástrofes e incidentes sobre infraestruturas estratégicas é padronizado internacionalmente pelo Sistema de Comando de Incidentes (\textit{Incident Command System} - ICS). O protocolo operacional preconiza a elaboração regular do boletim de situação \textbf{ICS-209} (\textit{Incident Status Summary}), que compila a localização do evento, ativos sob ameaça imediata, condições meteorológicas locais, recursos mobilizados e estimativas de contenção.

No âmbito de DataOps e engenharia de software contemporânea, a evolução de painéis analíticos passivos para sistemas proativos de resposta apoia-se em \textbf{Agentes de Inteligência Artificial} orquestrados pelo protocolo aberto \textit{Model Context Protocol} (MCP). Diferentemente de interfaces de chat genéricas vulneráveis a alucinações numéricas, a arquitetura com servidores MCP expõe ferramentas determinísticas tipadas diretamente sobre bancos analíticos colunares (Apache Druid e DuckDB) e grafos de topologia de rede (NetworkX).

O agente recebe comandos em linguagem natural emitidos pelo operador do centro de controle, traduz a intenção para dialeto SQL analítico por meio de técnicas de \textit{Few-Shot Prompting} ancoradas em consultas canônicas (\textit{Golden Queries}), avalia a proximidade geográfica contra ativos cadastrados e sintetiza o boletim ICS-209 em formato PDF e Markdown em menos de 30 segundos, disparando alertas enriquecidos para canais operacionais (Telegram, webhooks HTTP e protocolo XML \textit{Common Alerting Protocol} - CAP v1.2 da Defesa Civil).
